\section{\textit{Hyperparameter Tuning}}\label{sec:hyperparameter_tuning}

Dalam membangun model \textit{machine learning}, terdapat dua jenis parameter, yakni parameter internal yang dipelajari secara otomatis dari data oleh algoritma selama proses pelatihan dan \textit{hyperparameter} yang konfigurasinya harus ditentukan secara manual oleh pengembang sebelum proses pelatihan dimulai (Géron, 2019). \textit{Hyperparameter tuning} adalah tahapan kritis di mana pengembang bereksperimen mencari kombinasi \textit{hyperparameter} yang paling optimal guna menghasilkan model dengan performa prediksi terbaik dan minim tingkat kesalahan generalisasi pada data baru.

Sebuah konsep yang menjadi tujuan utama dalam proses \textit{tuning} adalah mencapai apa yang sering diistilahkan sebagai \textit{Sweet Spot}. Menurut Kuhn dan Johnson (2013) dalam \textit{Applied Predictive Modeling}, \textit{Sweet Spot} merupakan titik keseimbangan kompleksitas ideal di mana sebuah algoritma tidak terjebak pada kondisi \textit{underfitting} (terlalu sederhana hingga gagal menangkap pola) maupun \textit{overfitting} (terlalu rumit hingga menghafal gangguan atau \textit{noise} pada data latih secara presisi namun gagal saat digunakan memprediksi data nyata). Menemukan \textit{Sweet Spot} memastikan bahwa model benar-benar mempelajari pola fundamental yang relevan.

Selain itu, pemilihan arsitektur model dalam proses optimasi hyperparameter sangat dipandu oleh prinsip \textit{Parsimony} (\textit{principle of parsimony} atau \textit{Occam's Razor}). Prinsip \textit{parsimony} menyatakan bahwa jika terdapat beberapa model dengan performa prediksi yang serupa, maka model dengan struktur atau jumlah parameter yang paling sederhana dan paling transparan harus selalu diutamakan (Kuhn \& Johnson, 2013). Model yang menjunjung tinggi \textit{parsimony} jauh lebih efisien dalam hal komputasi, berjalan lebih cepat saat diimplementasikan ke dalam program nyata (\textit{deployment}), serta jauh lebih mudah dijelaskan alur cara kerjanya (\textit{interpretability}) kepada pemangku kepentingan tanpa mengorbankan akurasi yang berarti.

Salah satu metode modern yang terbukti sangat andal dan umum digunakan untuk mencari \textit{Sweet Spot} tersebut secara komputasional tanpa menguji keseluruhan probabilitas adalah \textit{Randomized Search Cross-Validation} (\texttt{RandomizedSearchCV}). Berbeda dengan metode penelusuran tradisional yang bersifat komprehensif (\textit{Grid Search}) sehingga memakan waktu ekstrem akibat permutasi kombinasi jutaan nilai, metode penelusuran acak menyelesaikan masalah tersebut dengan mengevaluasi sejumlah kombinasi acak dengan rentang iterasi dan probabilitas tertentu (Bergstra \& Bengio, 2012). Secara matematis dan probabilistik, mengambil sampel acak dalam jumlah yang memadai dari kumpulan \textit{hyperparameter} dapat menemukan performa sistem yang nyaris identik dengan hasil penelusuran \textit{grid} lengkap yang komprehensif. Pendekatan ini memungkinkan optimasi model dalam waktu yang dipercepat (Raschka \& Mirjalili, 2019).
